% TOP_PROBLEM: {"problemId": "problem.constant-parameter-quantum-locally-testable-codes", "canonicalTitle": "Constant-Parameter Quantum Locally Testable Codes", "releaseRank": 354, "primaryDomain": "quantum_information_computation", "primaryDomainLabel": "Quantum information & computation", "displayStatus": "open_with_solved_subcases", "releaseStatus": "open_with_solved_subcases"}
% !TeX program = lualatex
% ATTEMPT_STATUS: unresolved
\documentclass[11pt]{article}
\usepackage[margin=25mm]{geometry}
\usepackage{fontspec}
\setmainfont{Latin Modern Roman}
\usepackage{amsmath,amssymb,mathtools}
\usepackage{unicode-math}
\setmathfont{Latin Modern Math}
\usepackage{xurl,hyperref}
\hypersetup{colorlinks=true,urlcolor=blue}
\setcounter{secnumdepth}{0}
\setlength{\parindent}{0pt}
\setlength{\parskip}{5pt}
\setlength{\emergencystretch}{3em}
\begin{document}
\section*{\texorpdfstring{Constant-Parameter Quantum Locally Testable Codes}{Constant-Parameter Quantum Locally Testable Codes}}\label{354}
\subsection{Definitions and mathematical statement}
In the qubit convention, let $C\subset({\mathbb{C}}^2)^{\otimes n}$ be the common kernel of commuting projectors $\Pi_1,\ldots,\Pi_m$, each supported on at most $w$ qubits, and $H=\sum_i\Pi_i$. Let $C_t$ be the span of states obtained from $C$ by operators supported on at most $t$ qubits, $P_t$ its orthogonal projector, and
\[
 D_C=\sum_{t=1}^n t(P_t-P_{t-1}).
\]
Soundness $\rho$ means $H/m\succeq(\rho/n)D_C$, with $\succeq$ denoting positive-semidefinite order. Rate is $k/n$, where $k=\log_2\dim C$, and code distance is the least support size of an operator whose compression to $C$ is not scalar. The target is a family with $n\to\infty$, $k/n\ge r_0>0$, distance $d/n\ge\delta_0>0$, $\rho\ge\rho_0>0$, and $w\le w_0$, all constants independent of $n$. The source's fixed-local-dimension variant uses the corresponding logarithmic rate. Constant average check weight alone is insufficient.
\par\smallskip\textit{Formulation and concept references:} \href{https://errorcorrectionzoo.org/c/qltc}{[S1]}, \href{https://arxiv.org/html/2209.11405v3}{[E1]}.

\subsection{Short English statement}
Can quantum locally testable codes simultaneously have positive constant rate, relative distance, and soundness while every test acts on only a constant number of qubits?
\subsection{Research attempt}
\textbf{Status: unresolved by this attempt; a new external solution claim must be distinguished from this status.} On 17 September 2026, William Gay and Fernando Granha Jeronimo posted arXiv:2609.20780v1. Its Definition 1.2.1 uses the normalized Hamiltonian and quantum distance observable, and Theorem 2.4.3 asserts explicit binary CSS families of unbounded length with constant rate, relative distance, soundness, check weight, and qubit incidence [E2]. This matches the selected target, including the qubit and maximum-weight requirements. The inspected arXiv history on 27 September lists only v1, with no withdrawal or correction notice. This is a recent external claim, not an independently certified construction in this document; the inherited catalogue status has been preserved rather than used as evidence that the literature frontier is unchanged.

Two nearby September papers have different logical scopes. Bafna--Li--Nguyen [E3] assume a variant of a Reed--Solomon product-expansion conjecture. Cohen--Leigh--Reiner--Ta-Shma--Tzalik [E4] obtain a construction conditional on suitable lossless cubical complexes and explicitly do not construct those complexes. Their titles alone do not give additional unconditional proofs. Below I derive an exact stabilizer-code version of the selected operator inequality, construct informative families, and identify which simultaneous estimates an independent verification of an all-constant construction must establish.

\subsubsection{1. Exact reduction of the distance observable to syndrome lengths}
Let $g_1,\ldots,g_m$ be commuting Hermitian $n$-qubit Pauli operators, with nonzero simultaneous $+1$ eigenspace $C$. Thus the group $S$ that they generate does not contain $-I$. Put
\[
 \Pi_i=(I-g_i)/2,\qquad H=\sum_i\Pi_i.
\]
The generators need not be independent. If $S$ has binary rank $r$, then $\dim C=2^{n-r}$. A \emph{valid syndrome} is a vector $s\in\mathbb F_2^m$ for which the simultaneous eigenspace
\[
 E_s=\{v:g_iv=(-1)^{s_i}v\text{ for all }i\}
\]
is nonzero. Relations among the listed $g_i$ impose the corresponding linear relations on $s$. Let $Q_s$ be the orthogonal projection onto $E_s$.

Ignoring scalar phases, a Pauli is represented by $(x,z)\in\mathbb F_2^{2n}$ and has support size
\[
 |(x,z)|=|\operatorname{supp}x\cup\operatorname{supp}z|.
\]
Its syndrome records its commutation signs with the $g_i$. Define
\[
 \ell(s)=\min\{\operatorname{wt}(P):P\text{ a Pauli with syndrome }s\}.
 \tag{354.1}
\]
Every valid syndrome has such a Pauli. To justify this without an assumption of independent listed checks, identify $S$ modulo phases with an isotropic subspace $W$ of the nondegenerate symplectic space $\mathbb F_2^{2n}$. Restriction of the symplectic pairing gives a surjection onto $W^*$: a linear functional on $W$ extends to the whole space, and nondegeneracy represents it by a vector. This supplies precisely the allowed commutation patterns. A Pauli with syndrome $s$ maps $C$ isomorphically onto $E_s$, so each $E_s$ has dimension $2^{n-r}$.

\textbf{Lemma 354.1.} For this stabilizer presentation,
\[
 C_t=\bigoplus_{\ell(s)\le t}E_s,\qquad
 D_C=\sum_s\ell(s)Q_s,\qquad
 H=\sum_s|s|Q_s.
 \tag{354.2}
\]
Consequently the best soundness constant for this particular list of tests is
\[
 \rho_{\max}=\min_{\substack{s\ne0\\s\ \text{valid}}}
                 \frac{n|s|}{m\ell(s)},
 \tag{354.3}
\]
provided there is a nonzero valid syndrome.

\emph{Proof.} Any operator supported on a fixed set of at most $t$ qubits is a linear combination of Paulis supported there. Its image of $C$ lies in the sum on the right of the first identity. Conversely, for each syndrome of minimum weight at most $t$, a minimizing Pauli maps all of $C$ onto $E_s$. The two spans are therefore equal. Their orthogonal projectors increase by exactly the syndrome spaces with $\ell(s)=t$. The definition of $D_C$ gives the second identity. On $E_s$, the check $\Pi_i$ acts by the scalar $s_i$, proving the third. Both operators are diagonal on this orthogonal decomposition; positive-semidefinite comparison is exactly scalar comparison on each $E_s$. This proves (354.3). $\square$

This is also a reduction for arbitrary superpositions and mixed states. If a unit vector decomposes as $v=\sum_sv_s$, then
\[
 \langle v,Hv\rangle=\sum_s|s|\|v_s\|^2,
 \qquad
 \langle v,D_Cv\rangle=\sum_s\ell(s)\|v_s\|^2.
\]
No assumption that the tested state itself is a single Pauli error is needed. Conversely, a violating syndrome subspace already contains a pure-state counterexample.

\subsubsection{2. Local testability and logical distance control different kernels}
The minimum in (354.1) is taken over errors with a prescribed syndrome, including changes by all zero-syndrome Paulis. This is not the same minimum as the code distance. For $k=n-r>0$, the latter is
\[
 d=\min\{\operatorname{wt}(P):P\text{ commutes with }S,
                       \ P\notin S\text{ modulo scalar phases}\}.
 \tag{354.4}
\]
Here is a verification against the definition in the statement. A Pauli with nonzero syndrome maps $C$ to an orthogonal eigenspace and has zero compression. A Pauli in $S$ has scalar compression. A zero-syndrome Pauli outside $S$ has nonscalar compression: in the symplectic quotient $W^\perp/W$, it has a symplectic partner represented by another zero-syndrome Pauli. The two induced invertible operators on $C$ anticommute, which is impossible if either is scalar. Finally, expanding an arbitrary operator supported on $t$ qubits in the Pauli basis shows that if its compression is nonscalar, some Pauli of weight at most $t$ already has nonscalar compression. These observations prove (354.4).

In particular, a large-weight logical operator still has syndrome zero and sends $C$ to itself. Its contribution to $D_C$ is zero. It would be incorrect to substitute the weight of a chosen logical representative for $\ell(0)=0$ in the testing inequality. A soundness estimate concerns repair to \emph{some} code state, while distance protects the distinctions between states inside the code.

\subsubsection{3. The CSS reduction with the normalization constants retained}
Let $A\in\mathbb F_2^{m_X\times n}$ and $B\in\mathbb F_2^{m_Z\times n}$ satisfy
\[
 AB^T=0,
\]
and take $X$-checks from the rows of $A$ and $Z$-checks from the rows of $B$. Assume $m_X,m_Z>0$ and put $m=m_X+m_Z$. The rank of the stabilizer is $\operatorname{rank}A+\operatorname{rank}B$, hence
\[
 k=n-\operatorname{rank}A-\operatorname{rank}B.
 \tag{354.5}
\]
The Pauli $X^xZ^z$ has syndrome $(Az,Bx)$. Set
\[
 d_Z(z)=\operatorname{dist}(z,\ker A),\qquad
 d_X(x)=\operatorname{dist}(x,\ker B).
\]
For $s=(Az,Bx)$ one has
\[
 \max\{d_Z(z),d_X(x)\}\le\ell(s)
                           \le d_Z(z)+d_X(x).
 \tag{354.6}
\]
Indeed, each feasible Pauli representative has components in the prescribed affine cosets, so its union of supports has size at least each separate minimum. Choosing componentwise minimizers gives the upper bound. In a pure syndrome, equality is exact:
\[
 \ell(Az,0)=d_Z(z),\qquad \ell(0,Bx)=d_X(x).
 \tag{354.7}
\]
For example, in the first identity a pure $Z$ representative attains the lower bound.

Suppose the two classical parity-check maps satisfy, for all inputs,
\[
 \frac{|Az|}{m_X}\ge\frac{\alpha_Z}{n}d_Z(z),
 \qquad
 \frac{|Bx|}{m_Z}\ge\frac{\alpha_X}{n}d_X(x).
 \tag{354.8}
\]
Adding these inequalities and using (354.6) gives quantum soundness at least
\[
 \rho=\min\left\{\frac{m_X}{m}\alpha_Z,
                         \frac{m_Z}{m}\alpha_X\right\}.
 \tag{354.9}
\]
Conversely, quantum soundness $\rho$ implies, by the pure syndromes in (354.7), that the individual classical inequalities hold with
\[
 \alpha_Z=\rho\frac m{m_X},\qquad
 \alpha_X=\rho\frac m{m_Z}.
 \tag{354.10}
\]
These bounds suffice for constant-parameter comparisons when the check counts are comparable. They deliberately keep the weights of the two sampling distributions visible. Declaring the classical and quantum constants identical while normalizing each family separately would conceal a factor.

There are additional constraints beyond (354.8). The orthogonality $AB^T=0$ must hold; the ranks must leave $k$ linear in $n$; and the nontrivial classes in $\ker B/\operatorname{row}A$ and $\ker A/\operatorname{row}B$ must have linear minimum weight. Two unrelated good classical testers do not automatically satisfy these requirements.

\subsubsection{4. A constant-rate, constant-soundness family with distance one}
Fix $1\le a<n$ and let
\[
 C=|0\rangle^{\otimes a}\otimes(\mathbb C^2)^{\otimes(n-a)},
 \qquad g_i=Z_i\quad(1\le i\le a).
\]
There are $m=a$ checks, each of weight one, and $k=n-a$. Every syndrome $s\in\mathbb F_2^a$ can be corrected by flipping exactly its nonzero coordinates. No error on fewer coordinates can produce that syndrome, so $\ell(s)=|s|$. Formula (354.3) yields
\[
 \rho_{\max}=n/a.
 \tag{354.11}
\]
The convention in the question does not restrict soundness to be at most one; if desired one can simply report the weaker constant $\rho=1$. Taking $a=\lfloor n/2\rfloor$ gives positive constant rate, soundness, and locality. But a Pauli on any free qubit acts nonscalarly on $C$, so $d=1$.

The same obstruction prevents a direct use of a classical code as a quantum code. Let $L\le\mathbb F_2^n$ have positive dimension, and take the span of its computational basis vectors $|u\rangle$. Some coordinate $i$ is nonzero on some element of $L$. Since $0\in L$, the operator $Z_i$ has eigenvalue $+1$ on $|0\rangle$ and $-1$ on that codeword. Its compression is nonscalar, so the quantum distance is one, regardless of the classical minimum distance or classical local testability of $L$. Phase errors cannot be omitted from the selected target.

\subsubsection{5. The same code with very different checking Hamiltonians}
For $n\ge2$ consider the repetition subspace
\[
 C=\operatorname{span}\{|0^n\rangle,|1^n\rangle\}.
\]
First test its $n-1$ adjacent constraints $Z_iZ_{i+1}$. A bit-flip string $x$ and its complement have the same syndrome. These are the only two strings with that syndrome, so
\[
 \ell(s(x))=\min\{|x|,n-|x|\}.
 \tag{354.12}
\]
Extra $Z$ components of a Pauli do not help to produce a syndrome and cannot decrease the union of supports. Thus the formula also minimizes over all Paulis.

Flip an initial interval of length $t=\lfloor n/2\rfloor$. Exactly one adjacent constraint fails. Therefore
\[
 \rho_{\max}\le\frac{n}{(n-1)\lfloor n/2\rfloor}.
 \tag{354.13}
\]
In fact this is equality: every nonzero syndrome has at least one violated check, while its repair distance is at most $\lfloor n/2\rfloor$. Hence the path presentation has soundness of order $1/n$, despite constant check weight.

Now use every pair check $Z_iZ_j$, $i<j$, defining exactly the same code. There are $m=\binom n2$ checks of weight two. If $x$ has weight $t$, exactly $t(n-t)$ pairs cross its support. The syndrome ratio becomes
\[
 \frac{n\,t(n-t)}{\binom n2\min\{t,n-t\}}
       =\frac{2\max\{t,n-t\}}{n-1}.
\]
Consequently
\[
 \rho_{\max}=\frac{2\lceil n/2\rceil}{n-1}\ge1.
 \tag{354.14}
\]
Each qubit now occurs in $n-1$ checks. The selected formulation bounds each check's weight, not this incidence, so (354.14) is a legitimate constant-soundness example under its literal locality requirement. Nevertheless $k=1$ and $d=1$: $Z_1$ distinguishes the two basis states. It solves neither of the missing rate and relative-distance requirements.

This pair of examples also shows why testability belongs to a code together with its test distribution. Repeating \emph{every} check the same number of times leaves $H/m$ unchanged. Adding selected redundant checks can change the normalized Hamiltonian and its soundness substantially, even though its kernel is unchanged. One cannot treat redundancy as either always useful or always irrelevant.

\subsubsection{6. A finite verification that separates all four parameters}
For a stabilizer code on a fixed number of qubits, every quantity above can be calculated exactly over $\mathbb F_2$. Enumerate the $4^n$ Pauli vectors, compute each syndrome, retain the minimum union-support weight per syndrome, and separately find the minimum weight in $W^\perp\setminus W$. This uses exponential time and is only a finite diagnostic, not an efficient construction for increasing $n$.

For example, the four-qubit checks $XXXX$ and $ZZZZ$ commute and are independent. Thus $k=2$. All three nonzero syndromes are produced by one-qubit Paulis. The minimum nonzero syndrome weight is one, giving $\rho_{\max}=4/2=2$. There is no one-qubit zero-syndrome Pauli other than the identity; $X_1X_2$ is a zero-syndrome Pauli outside the stabilizer. Hence $d=2$. This provides the exact finite tuple
\[
 (n,k,d,m,w,\rho_{\max})=(4,2,2,2,4,2).
\]
It does not extend to an all-constant family merely by taking tensor products: disjoint copies leave $d=2$ while their total length increases.

Indeed, if codes on disjoint blocks have positive encoded dimension and minimum distance $d_0$, a minimum-weight logical operator on one block tensored with identities acts nonscalarly on the product code. The product distance is at most $d_0$; for identical blocks it is exactly $d_0$, since a Pauli of smaller total weight has scalar compression on every block. Thus independent replication preserves rate and locality but destroys relative distance. The same observation blocks the most direct amplification of any fixed good-looking finite example.

The accompanying exact diagnostic also checks the five-qubit code with generators $XZZXI$, $IXZZX$, $XIXZZ$, and $ZXIXZ$. It enumerates its stabilizer and syndrome spaces rather than presuming the advertised parameters, and checks (354.3), the two repetition presentations, and the ancillary-qubit family. These checks validate the reductions and examples; they do not validate the long construction in [E2].

\subsubsection{7. The first missing estimate in this attempt}
An independent all-constant construction must simultaneously exhibit unbounded $n$, linear stabilizer codimension, linear weight for nontrivial logical classes, bounded maximum check weight, and the uniform bound
\[
 |s|\ge c\,\frac mn\ell(s)
 \quad\text{for every valid syndrome }s,
 \tag{354.15}
\]
with one $c>0$ independent of $n$. Expansion only for sufficiently small errors does not automatically prove (354.15) for syndromes whose shortest repair has macroscopic weight. A bound on average check weight does not bound the worst test. And a classical soundness proof supplies neither commutation nor both logical distances.

The independent work here stops before constructing a commuting family with all these simultaneous properties. The external September claim [E2] explicitly asserts such a family and therefore requires consideration when updating the catalogue; [E3] and [E4] retain different unproved inputs. The local unresolved status records the scope of this attempt, not a rejection of [E2] and not an assertion that no solution has been announced.

\subsection{Sources}
\begin{itemize}
\item[S1]Error Correction Zoo, Quantum locally testable code (QLTC), maintained through 2026, accessed 2026-08-22.
\url{https://errorcorrectionzoo.org/c/qltc}.
\textit{Formulation source.}
\item[E1]Andrew Cross, Zhiyang He, Anand Natarajan, Mario Szegedy, and Guanyu Zhu, Quantum Locally Testable Code with Constant Soundness, arXiv:2209.11405v3 (2024), Section 2, Definition 2.4.
\url{https://arxiv.org/html/2209.11405v3}.
\textit{Added primary source: Quantum code fattening, the distance operator, and normalized soundness; the reported constructions are not a solution of the all-constant-parameter target.. Consulted 24 September 2026.}
\item[E2]William Gay and Fernando Granha Jeronimo, \emph{Asymptotically Good Quantum Locally Testable Codes}, arXiv:2609.20780v1, 17 September 2026.
\url{https://arxiv.org/html/2609.20780v1}.
\textit{Added primary source. Definition 1.2.1 and Theorem 2.4.3 inspected 27 September 2026; recent claim matching the exact all-constant binary CSS target. The full proof has not been independently certified here.}
\item[E3]Mitali Bafna, Anqi Li, and Quynh T. Nguyen, \emph{Good Quantum Locally Testable Codes from Product Expansion}, arXiv:2609.26735v1, 22 September 2026.
\url{https://arxiv.org/abs/2609.26735}.
\textit{Added primary source; abstract inspected 27 September 2026. The result assumes a variant of a product-expansion conjecture.}
\item[E4]Itay Cohen, Itai Leigh, Assaf Reiner, Amnon Ta-Shma, and Elad Tzalik, \emph{Good Quantum Locally Testable Codes from Lossless Cubical Complexes}, ECCC TR26-202, 21 September 2026.
\url{https://eccc.weizmann.ac.il/report/2026/202/}.
\textit{Added primary source; author abstract inspected 27 September 2026. Suitable lossless complexes are an input, not a construction supplied by that report.}
\end{itemize}
\end{document}
